In Group Relative Policy Optimization (GRPO)~\citep{shao2024deepseekmathpushinglimitsmathematical}, for each prompt $x$, we sample a group of $K$ trajectories $\mathcal{G}(x) = \{\tau^{(1)}, \dots, \tau^{(K)}\}$
% \[
%     \mathcal{G}(x) = \{\tau^{(1)}, \dots, \tau^{(K)}\},
% \]
where each trajectory $\tau^{(k)}$ is generated by the policy $\pi_\theta$ and receives a trajectory-level scalar return $R(\tau^{(k)})$.
GRPO constructs a \emph{relative} advantage within each group by subtracting a group-wise baseline from the individual return. Concretely, we define $\tilde{A}^{(k)} \;=\; R(\tau^{(k)}) - b(\mathcal{G}(x))$
% \begin{equation*}
%     \tilde{A}^{(k)} \;=\; R(\tau^{(k)}) - b(\mathcal{G}(x)),
% \end{equation*}
where $b(\mathcal{G}(x))$ is a baseline that depends only on the group statistics: $b(\mathcal{G}(x)) = \frac{1}{K} \sum_{j=1}^K R(\tau^{(j)})$.
% \begin{equation*}
%     b(\mathcal{G}(x)) = \frac{1}{K} \sum_{j=1}^K R(\tau^{(j)}).
% \end{equation*}
% This group-relative construction encourages trajectories that are better than their peers and discourages worse ones, while reducing variance without requiring an explicit learned value function.
The resulting policy gradient objective over a batch of groups is
\begin{equation*}
    J(\theta) \;=\; 
    \mathbb{E}_{x,\,\mathcal{G}(x)}\!\left[
        \frac{1}{K} \sum_{k=1}^K 
        \tilde{A}^{(k)} \sum_{t} \log \pi_\theta\!\big(a_t^{(k)} \mid s_t^{(k)}\big)
    \right],
\end{equation*}
where the inner sum runs over all decision tokens in the trajectory. In practice, we use a PPO-style clipped surrogate to stabilize training. Let $r_t^{(k)} = \frac{\pi_\theta(a_t^{(k)} \mid s_t^{(k)})}{\pi_{\theta_{\text{old}}}(a_t^{(k)} \mid s_t^{(k)})}$ denote the likelihood ratio; the GRPO loss is
\begin{equation}
\begin{split}
    \mathcal{L}_{\text{GRPO}}(\theta) &= - \mathbb{E}\Bigg[ \frac{1}{K} \sum_{k,t} \min\bigg( r_t^{(k)} \tilde{A}^{(k)}, \\
    &\quad \operatorname{clip}\big(r_t^{(k)}, 1-\epsilon, 1+\epsilon\big)\,\tilde{A}^{(k)} \bigg) \Bigg] \\
    &\quad + \beta\,\mathbb{E}\left[ D_{\mathrm{KL}}\big(\pi_\theta \,\|\, \pi_{\text{ref}}\big) \right]
\end{split}
\label{eq:grpo_loss}
\end{equation}
where $\epsilon$ controls the clipping range, $\beta$ weights a KL penalty to a reference policy $\pi_{\text{ref}}$, and the expectation is taken over prompts, groups, and time steps.

% \subsection{Memory Management}

% LLM agents take the entire interactions as their inputs (states) to make decisions, which becomes cumbersome (and sometimes infeasible) as the context length grows beyond the model’s limit. In practice, we need a proper memory management in order to increase the efficiency. The goal is to allow the LLM to operate on a summarized state rather than the full history. Formally, let the interaction history be:
% \begin{equation}
%     h=O_0,A_0,O_1,A_1, \dots, O_{t-1}, A_{t-1}, O_{t}
% \end{equation}
% The LLM agent instead observes a summarized state $z=f(h)$, where $f$ represents any kinds of summarization methods that compress the long interaction history into a compact state. Notably, the overall agent process satisfies Markov property only when $f$ captures the sufficient statistics of $h$. In other words, any future outcomes can be sufficiently predicted with $z$:
% \begin{equation}
%     p(\tau | z)=p(\tau | h)
% \end{equation}
% where $p(\tau | h)=p(O_{t+1},O_{t+2},O_{t+3},\dots | h, A_t, A_{t+1}, A_{t+2}, \dots)$.